\setcounter{framenumber}{0}
%21

\begin{frame}
  \titlepage
\end{frame}

\begin{frame}{Learning Goals}
By the end of this lecture, students should be able to:
\begin{itemize}
    \item define and interpret PDFs and CDFs;
    \item compute probabilities by integrating a PDF;
    \item recognize valid PDFs;
    \item compute expectation for continuous random variables;
    \item define and use the Uniform distribution;
    \item understand the universality of the Uniform distribution.
\end{itemize}
\end{frame}

\lecturepart{Continuous Random Variables}

\begin{frame}{From Discrete to Continuous Random Variables}
Discrete random variables take values in a finite or countably infinite set.

\vspace{0.25cm}

Continuous random variables can take any value in an interval, such as
\[
(0,1),\qquad (a,b),\qquad [0,\infty),\qquad \mathbb{R}.
\]

\vspace{0.25cm}

\begin{keybox}{Main difference}
For a continuous random variable,
\[
\Prob(X=x)=0
\]
for every individual point $x$.
\end{keybox}
\end{frame}

\begin{frame}{Probability Density Functions}
\begin{keybox}{Definition}
For a continuous random variable $X$, the probability density function, or PDF, is a function $f$ such that
\[
\Prob(X\in A)=\int_A f(x)\,dx, \quad \forall A\subseteq \mathbb{R}
\]
\end{keybox}

\vspace{0.25cm}

The value $f(x)$ itself is not a probability (could be greater than 1).

\vspace{0.25cm}

Probabilities are areas under the density curve.
\end{frame}

\begin{frame}{How to verify a function is a PDF?}
A function $f$ is a valid PDF if it satisfies two properties, namely

\[
f(x)\geq 0
\qquad \text{for all }x,
\]

and

\[
\int_{-\infty}^{\infty} f(x)\,dx=1.
\]

\vspace{0.25cm}

\begin{keybox}{Analogy}
For discrete random variables:
\[
p_X(x)\geq 0, \quad \sum_x p_X(x)=1.
\]

For continuous random variables:
\[
f_X(x)\geq 0, \quad \int_{-\infty}^{\infty} f(x)\,dx=1.
\]
\end{keybox}
\end{frame}

\begin{frame}{CDF and PDF}
CDF has the same formula for continuous and discrete r.v.s, i.e.
\[
F(x)=\Prob(X\leq x).
\]

For a continuous random variable with PDF $f$, according to the definition of PDF,

\[
F(x)=\int_{-\infty}^{x} f(t)\,dt.
\]

Note that , according to the definition of CDF,
\begin{equation*}
    \mathbb{P}(a< X\leq b) = F(b) - F(a) 
\end{equation*}
This, together with the Fundamental Theorem of Calculus, implies that
\begin{equation}
    \frac{d}{dx}F(x) = F^{\prime}(x) = f(x)
\end{equation}



For \textbf{continuous} r.v.s, since $\mathbb{P}(X=a)=0$ for any $a\in \mathbb{R}$,
\[
\begin{aligned}
\Prob(X\le a) &= \Prob(X<a),\\
\Prob(X\ge a) &= \Prob(X>a),\\
\Prob(a\le X\le b) &= \Prob(a<X\le b)
= \Prob(a\le X<b)
= \Prob(a<X<b).
\end{aligned}
\]








\end{frame}


\begin{frame}{What Does the Density \(f(x)\) Mean?}

For a continuous random variable,
\[
F(x)=\Prob(X\le x)
=
\int_{-\infty}^x f(t)\,dt.
\]

\vspace{0.2cm}

For a small number \(\epsilon>0\),
\[
\Prob(x\le X\le x+\epsilon)
=
F(x+\epsilon)-F(x).
\]

Therefore,
\[
\frac{\Prob(x\le X\le x+\epsilon)}{\epsilon}
=
\frac{F(x+\epsilon)-F(x)}{\epsilon}.
\]

Taking the limit as \(\epsilon\to0\),
\[
\lim_{\epsilon\to0}
\frac{\Prob(x\le X\le x+\epsilon)}{\epsilon}
=
\lim_{\epsilon\to0}
\frac{F(x+\epsilon)-F(x)}{\epsilon}.
\]

By the definition of derivative,
\[
f(x)=F'(x)
=
\lim_{\epsilon\to0}
\frac{F(x+\epsilon)-F(x)}{\epsilon}.
\]

Thus,
\[
{
f(x)
=
\lim_{\epsilon\to0}
\frac{\Prob(x\le X\le x+\epsilon)}{\epsilon}
}
\]

\end{frame}













\begin{frame}{Expectation for Continuous Random Variables}
\begin{keybox}{Definition}
If $X$ is continuous with PDF $f$, then
\[
\boxed{
\mathbb{E}(X)=\int_{-\infty}^{\infty} x f(x)\,dx.
}
\]
\end{keybox}

\vspace{0.25cm}

If the support is an interval, we integrate over that interval.

\vspace{0.25cm}

Expectation is still the balancing point of the distribution.
\end{frame}

\begin{frame}{Continuous LOTUS}
\begin{keybox}{LOTUS}
If \(X\) has PDF \(f\) and \(g:\mathbb{R}\to\mathbb{R}\), then
\[
\boxed{
\mathbb{E}(g(X))
=
\int_{-\infty}^{\infty} g(x)f(x)\,dx.
}
\]
\end{keybox}

\vspace{0.25cm}

\textbf{Proof idea.}
For a small interval \([x,x+\Delta x]\),
\[
\Prob(x\le X\le x+\Delta x)
\approx f(x)\Delta x.
\]

On this interval, \(g(X)\) is approximately \(g(x)\). Hence,
\[
g(X)\Prob(x\le X\le x+\Delta x)
\approx
g(x)f(x)\Delta x.
\]

Adding over many small intervals,
\[
\mathbb{E}(g(X))
\approx
\sum_i g(x_i)f(x_i)\Delta x.
\]

Taking the limit as \(\Delta x\to 0\), the Riemann sum becomes
\[
\mathbb{E}(g(X))
=
\int_{-\infty}^{\infty} g(x)f(x)\,dx.
\]

\end{frame}












\lecturepart{Uniform Distribution}
\begin{frame}{Uniform Distribution}

\begin{keybox}{Definition}
A random variable $U$ has the Uniform distribution on $(a,b)$ if its PDF is
\[
f(x)=
\begin{cases}
\dfrac{1}{b-a}, & a<x<b,\\
0, & \text{otherwise}.
\end{cases}
\]

We write
\[
U\sim \operatorname{Unif}(a,b).
\]
\end{keybox}

\begin{keybox}{PDF Verification}
\[
f(x)=\frac{1}{b-a}>0
\]
since \(b>a\). And
\[
\int_{-\infty}^{\infty} f(x)\,dx
=
\int_a^b \frac{1}{b-a}\,dx
=
\frac{1}{b-a}(b-a)
=
1.
\]
\end{keybox}

\end{frame}

\begin{frame}{Uniform CDF}
If
\[
U\sim \operatorname{Unif}(a,b),
\]
then its CDF is
\[
F(x)=
\begin{cases}
0, & x\leq a,\\[0.1cm]
\dfrac{x-a}{b-a}, & a<x<b,\\[0.2cm]
1, & x\geq b.
\end{cases}
\]

\vspace{0.25cm}

For the standard Uniform,
\[
U\sim \operatorname{Unif}(0,1),
\]
we have
\[
F(x)=x,\qquad 0<x<1.
\]
\end{frame}

\lecturepart{Universality of the Uniform}

\begin{frame}{Why the Uniform Distribution Is Special}
The Uniform distribution is a universal starting point for simulation.

\vspace{0.25cm}

Many algorithms first generate
\[
U\sim \operatorname{Unif}(0,1),
\]
then transform $U$ to obtain another distribution.

\vspace{0.25cm}

\begin{keybox}{Main idea}
A Uniform random variable can be transformed into a random variable with many other distributions.
\end{keybox}
\end{frame}

\begin{frame}{Inverse CDF Method}
Suppose $F$ is a continuous, strictly increasing CDF.

Let
\[
U\sim \operatorname{Unif}(0,1).
\]

Define
\[
X=F^{-1}(U).
\]

Then
\[
{
X \text{ has CDF } F.
}
\]

\vspace{0.25cm}

This is called the inverse CDF method or inverse transform sampling.
\end{frame}

\begin{frame}{Why the Inverse CDF Method Works}
Let
\[
X=F^{-1}(U).
\]

Then
\[
\Prob(X\leq x)
=
\Prob(F^{-1}(U)\leq x).
\]

Since $F$ is increasing,
\[
F^{-1}(U)\leq x
\quad \Longleftrightarrow \quad
U\leq F(x).
\]

Therefore,
\[
\Prob(X\leq x)
=
\Prob(U\leq F(x)).
\]

Because $U\sim\operatorname{Unif}(0,1)$,
\[
\Prob(U\leq F(x))=F(x).
\]

Thus $X$ has CDF $F$.
\end{frame}

\begin{frame}{Probability Integral Transform}
Conversely, if $X$ has continuous strictly increasing CDF $F$, then
\[
{
F(X)\sim \operatorname{Unif}(0,1).
}
\]

\vspace{0.25cm}

\begin{keybox}{Interpretation}
Applying the CDF to a random variable converts its value into a percentile.
\end{keybox}

\vspace{0.25cm}

For example, if exam scores have CDF $F$, then $F(X)$ is the percentile of a randomly selected student.
\end{frame}
\begin{frame}{Uniform Mean and Variance}
If $U\sim \operatorname{Unif}(a,b)$, then
\[
f_U(u)=\frac{1}{b-a},
\qquad a\le u\le b.
\]

\begin{keybox}{Mean}
\[
\mathbb{E}(U)
=
\int_a^b u\,\frac{1}{b-a}\,du
=
\frac{1}{b-a}\left[\frac{u^2}{2}\right]_a^b
=
\frac{b^2-a^2}{2(b-a)}
=
{\frac{a+b}{2}}.
\]
\end{keybox}

\begin{keybox}{Variance}
First,
\[
\mathbb{E}(U^2)
=
\int_a^b u^2\,\frac{1}{b-a}\,du
=
\frac{b^3-a^3}{3(b-a)}
=
\frac{a^2+ab+b^2}{3}.
\]
Thus,
\[
\operatorname{Var}(U)
=
\mathbb{E}(U^2)-[\mathbb{E}(U)]^2
=
\frac{a^2+ab+b^2}{3}
-
\left(\frac{a+b}{2}\right)^2
=
{\frac{(b-a)^2}{12}}.
\]
\end{keybox}
\end{frame}

\begin{frame}{Location-Scale Transformation}
If
\[
U\sim \operatorname{Unif}(0,1),
\]
then
\[
a+(b-a)U\sim \operatorname{Unif}(a,b).
\]

\vspace{0.25cm}

This transforms:
\[
(0,1)
\quad \longrightarrow \quad
(a,b).
\]

\vspace{0.25cm}

\begin{keybox}{Useful idea}
Start with the standard Uniform and shift/scale to get any Uniform interval.
\end{keybox}
\end{frame}

\begin{frame}{Summary}
\begin{itemize}
    \item Continuous random variables have probabilities over intervals, not points.

    \item A PDF integrates to probabilities:
    \[
    \Prob(a<X<b)=\int_a^b f(x)\,dx.
    \]

    \item A valid PDF is nonnegative and integrates to $1$.

    \item The CDF is accumulated area:
    \[
    F(x)=\int_{-\infty}^x f(t)\,dt.
    \]

    \item Uniform probabilities are proportional to interval length.

    \item The Uniform distribution is universal through the inverse CDF method.
\end{itemize}

\vspace{0.25cm}

\begin{keybox}{Next lecture}
Normal and Exponential distributions.
\end{keybox}
\end{frame}